The current work is important for safety-related transients, as the \gls{MSR} specific \gls{DNP} group parameters will impact the results of these transients.
Depending on the properties of the specific reactor, which affect the group parameters, the static \gls{DNP} group parameter transient simulation may generate overly or insufficiently conservative results.
This work will use safety-related transients to determine the effects of the improved \gls{DNP} group parameters in the \gls{MSRE}, and under what conditions, if any, the previous \gls{DNP} groups would predict safe operation while the improved groups would predict unsafe conditions.

An important part of ensuring safe reactor operation is making sure safety-related transient events do not disrupt safe reactor operations.
These transients can be caused by \glspl{PIE}, which are defined in the \gls{IAEA} Safety Standard Series No. NS-R-1 \cite{international2000iaea}.
Of particular importance are \glspl{PIE} that potentially lead to \glspl{DBA}.
Although No. SSR-2/1 has superseded No. NS-R-1, the \glspl{PIE} and \glspl{DBA} for \glspl{MSR} have not changed \cite{international2016iaea}.
Although there are not specific events listed for \glspl{MSR}, guidance is provided on categories of events that are of interest..

These events include equipment failures, such as pumps failing or pipes rupturing; human error, such as improper maintenance or faulty operator actions; external events, such as earthquakes; and combinations of events.
For combinations of events, the \gls{IAEA} specifies that these events should make sense in combination and in the time frame modeled.
For example, if a particular event lasts a long time or has a long lasting effect, then that event could possibly overlap with another unrelated event.
Alternatively, dependent events could potentially occur at the same time.


Safety-related transients for \glspl{MSR} include unprotected loss of heat sink (ULOHS), unprotected loss of flow (ULOF), \gls{UTOP}, and \gls{FSOC} \cite{li2015transient, daher2022improvement}.
A \gls{ULOHS} transient is when the heat exchanger experiences a failure, possibly caused by a pump failure in the intermediate loop, causing a loss of cooling of the reactor.
This causes higher fuel salt temperatures, which can cause adverse chemical effects and increase thermal stress \cite{belles2020proposed}.
For the \gls{MSFR}, Brovchenko et al. showed that slow \gls{ULOHS} transients that take place on time scales longer than a minute minimize the increase in fuel salt temperature \cite{brovchenko2013design}.
A \gls{ULOF} transient is when the flow rate in the primary loop becomes zero, caused by pump failure in the primary loop.
The \gls{ULOF} transient can lead to an increase in the fuel salt temperature, similar to the \gls{ULOHS} transient.
A \gls{UTOP} transient is a reactivity insertion event, which could be caused by various means, such as removing poisons, adding fissile material, or temperature changes.
This transient also causes an increase in the fuel salt temperature, similar to the \gls{ULOHS} and \gls{ULOF} transients.
An \gls{FSOC} transient is when the fuel salt temperature decreases, which could be caused by an increased heat removal rate from the heat exchanger \cite{belles2020proposed}.
Overcooling the fuel salt can lead to flow blockages from the salt freezing as well as a reactivity insertion from the decrease in fuel salt temperature.


%These safety-related transients are important for the current work, as the tabulated \gls{DNP} group parameters can be used to more rapidly simulated these transients.
%Additionally, due to the coupling of the neutronics and thermal-hydraulics in calculation of the tabulated effective group parameters, the results are more accurate than they would be using more traditional transient analysis methods.
%The method used to simulate the transient is also an important consideration.